\section{Conclusion}

This thesis presents an analysis of AI persuasiveness in a debate platform, investigating whether
AI opponents effectively shift human participant beliefs compared to human-human debate.
The primary empirical finding is that AI debate produces \textit{null belief change}
with only a \textit{directional (non-significant) confidence increase}. This pattern is
methodologically informative,
offering implications for AI's role in persuasion, argumentation training, and human-AI collaboration.

\subsection{Primary Finding: The Confidence-Belief Dissociation}

Across four pre-specified confirmatory models with robust standard errors,
AI debate produced no significant effect on participant belief change relative to human-human debate
(all $p > 0.35$, $d = 0.286$, observed power $= 20.0\%$).
Confidence models showed a consistent directional increase that did not reach statistical significance
($\beta = 5.2$--$5.8$ points, 95\% CI crossing 0, $p = 0.17$--$0.21$).

Because baseline confidence was recorded as 0 for all participants, this should be interpreted as
between-condition differences in post-debate confidence levels rather than true baseline-to-post change.

This pattern is consistent with classical persuasion theory, which often models belief and
confidence as closely coupled constructs \cite{petty1986_elm,eagly1993_psychology_of_attitudes}.
In the AI debate context, participants apparently update their \textit{confidence in debate performance}
without updating their \textit{core stance conviction}. They finish debates feeling more certain
in direction, without reliable evidence of shifted underlying belief.

\subsection{Mechanism: Illusory Success, Not Defensive Conviction}

We hypothesized that the confidence gain reflected a ``defensive confidence'' mechanism:
participants who argued critically should feel more convinced of their position after ``winning'' debate exchanges.
However, correlation analysis of the qualitative subsample revealed the opposite:
participants who engaged in high pushback (critical argumentation) showed \textit{lower} confidence gain
($r = -0.121$, $p = 0.707$, 95\% CI $[-0.650,0.487]$), not higher. Jackknife resampling
showed a directional negative range ($r_{jackknife} \in [-0.372,-0.003]$), but with weak inferential certainty.

One plausible interpretation is \textit{illusory success}:
\begin{itemize}
    \item Passive participants perceived themselves as ``winning'' against a non-concessive opponent,
          inflating confidence without substance.
    \item Active arguers recognized the AI's rhetorical resilience and updated confidence downward.
\end{itemize}

This interpretation is consistent with AI design (the AI is prompted to defer politely but never concede major points)
and participant framing (belief measurements follow immediately after debate, anchoring confidence to the stated position).

\subsection{Theoretical Implications}

\subsubsection{AI as a Credible but Non-Persuasive Interlocutor}

Participants may treat AI as a serious intellectual opponent (evidenced by directional confidence differences and continued engagement),
but they do not grant the AI epistemic authority
to shift core beliefs. This suggests:

\begin{enumerate}
    \item Participants recognize AI as non-human and apply a discount to its authority on value-laden or identity-central topics.
    \item Confidence gains are divorced from belief shifts, perhaps because confidence is self-referential
          (``How sure am I of my position?'') while belief resistance reflects source skepticism (``How credible is the AI?'').
    \item For topics where AI might have genuine epistemic authority (factual matters, logical consistency),
          the debate platform may not capture belief shifts because participant stance choices are often taken on priors.
\end{enumerate}

\subsubsection{Implications for AI Pedagogy and Argumentation Training}

If AI debate directionally boosts confidence without shifting beliefs, its value may lie in \textit{self-efficacy training}
rather than \textit{persuasion}. Participants may benefit from:
\begin{itemize}
    \item Increased confidence in their own reasoning and argumentation skills,
    \item Reduced anxiety about defending positions,
    \item Practice formulating counterarguments in real time.
\end{itemize}

However, if the intent is to expose participants to genuinely compelling opposing views and shift deeply held beliefs,
AI debate alone is insufficient.

\subsection{Methodological Strengths}

This analysis demonstrates rigorous small-sample inference through:
\begin{itemize}
    \item \textbf{Pre-specified model sequence}: Four confirmatory models locked before data analysis.
    \item \textbf{Robust standard errors}: HC3 covariance matrices account for heteroscedasticity.
      \item \textbf{Inferential transparency}: confidence models are reported with explicit effect sizes,
                    confidence intervals, and corresponding $p$-values.
                                    \item \textbf{Transparency about underpowering}: With model sample $n = 89$, observed power for belief effect detection is $20.0\%$.
                                          The null result is stable across reported sensitivity checks but remains compatible with limited power.
    \item \textbf{Mechanism validation attempted}: Correlation test falsified initial hypothesis,
              demonstrating intellectual honesty rather than confirmatory bias.
    \item \textbf{Data quality audits}: Floor-ceiling effect checks and baseline confidence missing-data detection.
\end{itemize}

\pagebreak
\section{Limitations}

\subsection{Sample Size and Power}

With model sample $n = 89$, this study is underpowered for \textit{belief} effects ($d = 0.286$, required $n \approx 175$ per group for $80\%$ power).
The null belief finding is stable across robustness checks,
but replication with a larger sample is warranted.

\subsection{Measurement and Scale Issues}

\begin{itemize}
    \item \textbf{Belief scale coarseness}: 0--100 continuous scale may not capture nuanced within-debate shifts,
          particularly for topics with high pre-engagement ceiling effects.
    \item \textbf{Baseline confidence missing}: All participants recorded 0 confidence pre-debate.
          While this does not invalidate inference, it renders ``confidence change'' equivalent to ``post-debate confidence,''
          complicating interpretation.
    \item \textbf{Reflection quality proxy}: Text length is a crude proxy for engagement quality.
          Ideally, qualitative codes would be applied exhaustively rather than to a stratified sample.
\end{itemize}

\subsection{Sample Representation Bias}

Qualitative analysis was performed on a stratified sample ($n \approx 24$ reflections)
that overrepresents human-human debates ($50\%$ vs. $20.5\%$ in population).
AI conditions are each underrepresented in the qualitative sample ($16.7\%$ each)
relative to their population shares (firm $28.4\%$, balanced $27.4\%$, open-minded $23.7\%).
This bias limits the representativeness of mechanism findings and should be considered exploratory.

\subsection{Generalizability}

Participants were recruited online and self-selected into debate participation, possibly skewing toward
debate-curious, relatively confident individuals. Results may not generalize to the broader population.
Additionally, topics were participant-selected rather than experimenter-controlled, introducing unmeasured heterogeneity.

\subsection{AI Realism and Authenticity}

The AI opponent is a prompted language model, not a persistent agent with learning. The pattern of
"never conceding", while methodologically clean, may not reflect authentic disagreement dynamics
in real-world human-human or human-AI debate.

\pagebreak
\section{Future Work}

Future research directions address the dissociation finding and explore boundary conditions:

\subsection{Longitudinal Belief Tracking}

Does the confidence boost persist, and does it eventually translate into belief shifts?
Measuring participant beliefs days or weeks post-debate, with follow-up booster debates,
could reveal whether initial illusory confidence resolves into genuine persuasion with repeated exposure.

\subsection{Factual vs. Value-Based Topics}

The current sample spans diverse topics (political, moral, empirical). Stratified analysis by topic type
could test whether AI is more persuasive on factual matters (where it has genuine epistemic advantage)
than value-laden stances (where source credibility matters more).

\subsection{Belief Attribution Separately Randomized}

Implement a $2 \times 2$ design crossing actual opponent type (AI vs. human) with \textit{belief about} opponent type.
This would isolate whether confidence effects arise from real AI properties or from participants' \textit{belief}
that they debated an AI.

\subsection{Qualitative Deepening and Full Coding}

Apply thematic coding exhaustively to all reflections and arguments, not just a stratified sample.
Integrate quantitative outcome measures with qualitative content analysis to triangulate mechanism hypotheses.

\subsection{AI Personality Fine-Tuning}

Current AI personalities (firm, balanced, open-minded) show no effect on belief outcomes.
Experiment with AI personalities that explicitly concede strategic points, express uncertainty, or
adopt Socratic questioning. Would such modifications increase persuasiveness?

\subsection{AI as Knowledge Authority vs. Opinion Debater}

Separate debates into two tracks: (1) factual knowledge debates (e.g., scientific claim verification)
where AI has genuine advantage, and (2) value/opinion debates where human credibility dominates.
Hypothesize that AI persuasiveness is moderated by topic type.

\subsection{Practical Application: Argumentation Tutoring}

Since AI debate shows directional confidence increases, deploy the platform as a confidence-building tool in
argumentation curricula. Compare students' debate performance pre- and post-AI-debate practice.
Test whether confidence gains transfer to performance against human opponents.

\pagebreak
\section{Final Remarks}

This thesis contributes to our understanding of AI persuasiveness by demonstrating that
AI in debate format does not produce reliable pooled effects on either belief change or confidence change
in this pilot, while still supporting useful engagement and process diagnostics.
This pattern motivates a nuanced model of AI's role:
not as a replacement for human-led persuasion or education, but as a complement for practice,
confidence building, and intellectual sparring.

The methodological framework---locked designs, robust inference, and mechanism validation---
provides a template for rigorous empirical work on AI-human interaction in persuasive contexts.
As AI becomes increasingly integrated into learning and decision-support systems,
understanding its limitations and comparative advantages remains critically important.
Numeric claims in this chapter are synchronized to Appendix
Table~\ref{tab:headline_numbers}.
